% New increment only. No prior source is modified.
\documentclass[11pt]{article}
\usepackage[T1]{fontenc}\usepackage[utf8]{inputenc}\usepackage{lmodern}
\usepackage[margin=.9in]{geometry}
\usepackage{amsmath,amssymb,amsthm,mathtools,microtype,longtable,array,booktabs}
\usepackage[hidelinks]{hyperref}
\setlength{\emergencystretch}{3em}
\newtheorem{theorem}{Theorem}[section]
\newtheorem{proposition}[theorem]{Proposition}
\newtheorem{corollary}[theorem]{Corollary}
\newcommand{\dd}{\,d}
\newcommand{\K}{\mathcal K}
\newcommand{\E}{\mathbb E}
\newcommand{\one}{\mathbf1}
\newcommand{\kf}{\mathsf K}
\title{Original-flow B-core audit:\\the nonvanishing cluster-2 output-rotation channel}
\author{Standalone analytical increment for review}
\date{Version 1 --- 2 November 2026}
\begin{document}\maketitle
\begin{abstract}
Step 2 of the proposed S6 contract fails on the original initialized
deep-regime flow itself. With the checkpoint's unchanged definitions,
its remainder \(r_2\) contains cluster-2 output rotation of the strong
core. On WIT's fixed learned windows,
\[
 q^{-2}r_2=
 \frac{\rho}{2}(1-n_\rho)(\zeta-X_*)\kf(\rho X_*)
 +O(q^{1/20}),
\]
uniformly on the specified parameter rectangle. The displayed leading
term has an explicit positive, \(q\)-independent lower bound.
Consequently \(r_2=O(q^{2+\kappa})\) is false for every \(\kappa>0\).
The exact factorizations remain valid; the proposed budget does not.
This increment proves that obstruction with all background labels
retained, and stops at Step 2. It does not claim the full B-core
theorem or proceed to signed-lag and producer-attribution tasks.
\end{abstract}

\section{Scope, conventions and hypothesis inventory}
\label{an01bcv1:inventory}
Use the initialized balanced untied Gaussian population of RTA and FWC,
\(S_0=q^{10}\), with their already fixed constant order and sufficiently
small \(q>0\). Keep WIT's original compact rectangle
\(\mathcal P_*=K_\rho\times I_\zeta\), where
\[
 K_\rho\subset[13/20,7/10],\qquad
 I_\zeta\subset[87/50,7/4],
 \qquad J=[s_*-r,s_*+r].
\]
The same matched clock is
\[
 \tau=\mu_1^2t_{\rm phys},\qquad
 \tau_h=t_n+\lambda^{-1}\log\frac{1-n}{n},\qquad
 s=\tau-\tau_h,\qquad \lambda=\rho^2.
\]
Here \(n\) is the fixed entry level, whereas
\(n_\rho(s)=(1+e^{-\lambda s})^{-1}\) is the comparison mass.
Write \(d\mu=m\,d\lambda_0\), \(M=\mu_\K\), \(N=\mu_S\);
the core \(\K\), selected cohort \(S\), primary remainder \(P_q\),
secondary cohort \(\Sigma\) and outside set \(B\) are the fixed-label
sets in FWC. No new cohort is selected.

For a receiver write \(a=a_\theta=(\cos\theta,\sin\theta)\),
\(b=a_\psi\). Let \(e_\xi=\xi_q-f(\xi_q)\), with
\(\xi_q=(\sin q\zeta,-\cos q\zeta)\).
The checkpoint uses the \emph{unhalved} source matrix
\[
 R_p(\theta)=\E_{P_p}[(X-f(X))X^T\one_{\{a\cdot X>0\}}],
\]
and velocities
\[
 g_p=\tfrac12 b^TR_pa,\quad
 v_{\theta,p}=\tfrac12 b^TR_pa^\perp,\quad
 v_{\psi,p}=\tfrac12(b^\perp)^TR_pa.
\]
Thus \((\log m)'_p=2g_p\). RTA's matrix \(R^{(p)}\) equals
\(R_p/2\), and its full radial reaction \(h_p=2g_p\).
These conventions account for every mixture factor below.

\small
\begin{longtable}{>{\raggedright\arraybackslash}p{.22\textwidth}
 >{\raggedright\arraybackslash}p{.44\textwidth}
 >{\raggedright\arraybackslash}p{.25\textwidth}}
\toprule Source & Hypotheses and exact/leading scope & Disposition\\
\midrule\endhead
P \texttt{motion}
& Exact balanced characteristic flow, finite moment integrals, fixed
 probe; chain rule at gate-regular times, or selector/a.e. convention.
 All labels and all three motions are retained.
& Used. FWC supplies finite moments and RTA supplies the
 gate-separated core.\\
P \texttt{projection}
& Exact full output; \(P_1=N(e_1,q^2I)\), \(Q_1=1+q^2\);
 current probe-active/inactive partition; finite \(L^2(P_1)\) output.
 Uses \(\E_{P_1}X_1X_2=0\). No lag sign is assumed or concluded.
& Exact identity available. Not used to infer a signed lag.\\
P \texttt{factorizations}, \texttt{epsilonB}, \texttt{epsilonH}
& A measurable core \(C\) contained in the current active set,
 \(M_C>0\); normalized current mass averages, signed variables,
 finite integrals. No smallness of remainders is a hypothesis or
 conclusion of the algebraic identities.
& Use \(C=\K\). RTA gives strict activity, FWC gives \(M\ge1/2\).
 The proposed smallness is tested, not imported.\\
P \texttt{helpdefect}
& Same exact Gaussian setting and variables. The four-term identity
 is algebraic; its displayed Gaussian tail estimate additionally
 assumes \(a_1>0\) and finite \(\|e_2\|_{P_1}\).
& Hypotheses hold on the compact core chart, but its quantitative
 evaluation is not needed after the \(r_2\) obstruction.\\
P \texttt{GLC}, \texttt{GLCidentity}
& The preceding exact core definitions. A rate identity at
 chain-rule times; no autonomous closure or small-error assertion.
 Separate-sign criteria require the explicitly displayed margins.
& Exact bookkeeping retained. Its error budget cannot be replaced
 by \(o(q^2)\) under the proposed partition.\\
P \texttt{producer}
& An exhaustive output partition with integrable moments;
 same receiver gate in each term; same rate-weighted core integral.
 No evaluation at a mean receiver is allowed.
& Available exact identity; producer asymptotics not attempted
 after Step 2 fails.\\
P \texttt{tracking}
& Fixed initial-label subset and complement, Gaussian \(L^2\)
 projection, exact kernel evolution and regularity for differentiation.
 Positive damping estimate also needs \(|\theta|\le\theta_0<\pi/2\);
 the comparison bound assumes \(\kappa_A\ge\kappa_*>0\),
 \(|\eta_A|\le\eta_*\). It does not prove a positive lag.
& Not invoked. No forcing or signed-lag hypothesis is added.\\
P \texttt{lagtheorem}, \texttt{lag-pop}, \texttt{Xi-positive}
& Bounded-chart leading characteristic solution, \(M=1\),
 \(0\le N\le1\); positive measures; continuous interpretation at
 \(N=0\); finite coordinates for strict \(\Xi>0\); positive incoming
 lag for positivity preservation.
& Leading-only statements. Not substituted for exact original
 lag positivity; not needed for the obstruction.\\
P \texttt{atomiclag}, \texttt{atomicrates}
& Selected coherent leading branch, \(0<n_\rho<1\), \(M=1\);
 active strong probe for the displayed rate formula, inactive weak
 probe in its two-family representation.
& Only the leading rate identity is used as a consistency check.
 The original estimate is proved using RTA/FWC.\\
RTA \texttt{coregap}, \texttt{backgroundrate}, \texttt{ratetransfer}
& Original \(S_0=q^{10}\) family; sufficiently small q after
 fixed constants; original \(\mathcal P_*\times J\); matched clock.
 Selector source formulas, classical derivatives where defined.
& Unconditional inputs on exactly this window. Core gap
 \(\zeta-x\ge g_*/8\); supremum tracking \(O(q^{1/20})\).\\
FWC \texttt{corefield}, \texttt{distbound}, \texttt{background}
& Original deep-regime initialized family; fixed finite
 \(T>\max_\rho s_n\), bounded offsets, q last.
 Physical current masses/moments and unchanged cohorts.
& Choose fixed T covering J. Source-2 output velocity,
 reaction, masses and background estimates are available.\\
WIT \texttt{uniformtheorem}
& Selected leading branch and original rectangle/windows.
& Gives \(\zeta-X_*\ge g_*/4\) on J.
 Finite logistic time gives \(1-n_\rho>0\) uniformly.\\
WIT \texttt{regiontheorem}
& Selected branch, its ratio box, and specifically \(s\le0\).
& Not extended to an arbitrary point of J. No positivity of
 \(X_*\) is needed below.\\
\bottomrule
\end{longtable}\normalsize

\section{Exact factorization and the channel left in \texorpdfstring{\(r_2\)}{r2}}
\label{an01bcv1:channels}
Use the checkpoint's exact definitions on \(C=\K\):
\[
 z=a\cdot\xi_q>0,\quad v=-a^\perp\cdot\xi_q,\quad
 d=-e_\xi\cdot b^\perp,\quad c=e_\xi\cdot b,\quad
 T_2=R_2a^\perp,\quad j=b_1T_{2,1}.
\]
The scalar \(c\) here is the checkpoint's probe projection, not the
weak training response. Its sign is not imposed.
With bars denoting normalized current core averages,
\begin{align}
 B_0&=\tfrac M2Q_1\bar z\,\bar d\,r_{21},&
 H_{1,0}&=\tfrac M2\bar v\,\bar c\,\bar j,\notag\\
 D_1^\tau&=-B_0-\epsilon_B+r_1,&
 D_2^\tau&=H_{1,0}+\epsilon_H+H^{(2)}+r_2.
 \label{an01bcv1:exactfactor}
\end{align}
These remain exact. In particular
\begin{equation}
 H^{(1)}+H^{(2)}
 =\frac12\int_\K vc\,b^TR_2a^\perp\dd\mu
 =\int_\K vc\,v_{\theta,2}\dd\mu
 \label{an01bcv1:inputchannel}
\end{equation}
is the full core input-rotation channel, including both coordinates
of \(b^TR_2a^\perp\).

Define the full-residual core output and radial channels by
\[
 O_{2,\K}=\int_\K zd\,v_{\psi,2}\dd\mu,\qquad
 P_{2,\K}=-\int_\K 2g_2zc\dd\mu.
\]
P \texttt{motion} gives exactly
\begin{equation}
 r_2=O_{2,\K}+P_{2,\K}+D_{2,\K^c},
 \label{an01bcv1:r2identity}
\end{equation}
where \(D_{2,\K^c}\) retains all three motions of every non-core
label. Thus no algebraic convention permits one to call \(r_2\)
only a finite-noise error.

\begin{theorem}[Positive leading term in the exact \(r_2\)]
\label{an01bcv1:obstructiontheorem}
Put \(\chi=1/20\), \(L=1+\log(1/q)\),
\(\gamma_*=289169/924000\), and
\(\epsilon_{\rm bg}=q^8+L^2q^{\gamma_*}\).
There is a fixed C, independent of q, such that uniformly on
\(\mathcal P_*\times J\),
\begin{equation}
 \left|q^{-2}r_2-\mathcal O_*(\rho,\zeta,s)\right|
 \le C(q^\chi+q^{1/10}+q^2+\epsilon_{\rm bg})
 \le C_{\rm rem}q^\chi,
 \label{an01bcv1:r2limit}
\end{equation}
where the signed leading term is
\begin{equation}
 \mathcal O_*=\frac{\rho}{2}(1-n_\rho)
                 (\zeta-X_*)\kf(\rho X_*).
 \label{an01bcv1:Ostar}
\end{equation}
Let \(R\ge2\) bound the reached core and comparison charts and offsets
on a fixed FWC window covering J. Define
\begin{equation}
 b_J=|s_*|+r,\quad
 \delta_n=\frac1{1+\exp((49/100)b_J)},\quad
 c_{\rm out}=\frac{13g_*}{160}\delta_n\kf(-7R/10)>0.
 \label{an01bcv1:cout}
\end{equation}
Then \(\mathcal O_*\ge c_{\rm out}\). Consequently, after decreasing
q so that \(C_{\rm rem}q^\chi\le c_{\rm out}/2\),
\begin{equation}
 r_2(\tau_h+s,\xi_q)\ge\tfrac12 c_{\rm out}q^2
 \qquad ((\rho,\zeta,s)\in\mathcal P_*\times J).
 \label{an01bcv1:r2lower}
\end{equation}
All source rates here use the exact selector convention, with the
usual classical/a.e. interpretations; there is no core gate hit.
\end{theorem}

\begin{proof}
We first keep background in the precise observable needed, rather
than assuming it is absent. Set \(T_0=\Sigma\cup B\).
FWC and the RTA background proof give
\[
 \mu_{T_0}\le Cq^{\gamma_*},\qquad
 \mathcal Q_{T_0}\le Cq^2\epsilon_{\rm bg},\qquad
 |f_{T_0}(\xi_q)|\le\sqrt{\mu_{T_0}\mathcal Q_{T_0}}
                         \le Cq\epsilon_{\rm bg}.
 \tag{*}
\]
The last inequality follows from \(q^{\gamma_*}\le\epsilon_{\rm bg}\).
The bound before it follows directly by Cauchy--Schwarz applied
to the positive feature integral; \(\mathcal Q\) is FWC's squared
active-probe input integral. FWC gives exact primary inactivity,
so \(f=f_\K+f_{T_0}\) at this probe. RTA also gives
\[
 |D_2^\tau-D_2^{\K,\mathrm{self}}|\le Cq^2\epsilon_{\rm bg},
 \qquad
 D_2^{\K,\mathrm{self}}=-(\xi_q-f_\K)\cdot\dot f_{\K,2}.
 \tag{**}
\]
These bounds include all outside labels and probe-gate slices.

Replacing \(e_\xi\) by \(e_\xi^\K=\xi_q-f_\K\) in the core
input channel changes it by at most
\(|f_{T_0}|\,|\dot f_{\K,2}^{\,\theta}|\).
On the bounded core chart, FWC gives
\(v_{\theta,2}=O(q)\), so
\(|\dot f_{\K,2}^{\,\theta}|\le Cq\).
The change is therefore \(O(q^2\epsilon_{\rm bg})\).
Subtracting this input channel from (**) yields
\begin{equation}
 r_2=O_{2,\K}^{\rm self}+P_{2,\K}^{\rm self}
                       +O(q^2\epsilon_{\rm bg}).
 \label{an01bcv1:r2self}
\end{equation}
This is also a direct bound on the combined non-core and
full-output corrections in \eqref{an01bcv1:r2identity}; none is
discarded by ancestry or sign.

For compact core coordinates \(\theta=qx,\ \psi=qy\), put
\(A=\int_\K(\zeta-x)\dd\mu\), \(B_\zeta=\int_\K y(\zeta-x)\dd\mu\).
RTA's core expansion and the reached chart give
\[
 f_{\K,1}=qA+O(q^3),\quad f_{\K,2}=q^2B_\zeta+O(q^4),
 \quad z=q(\zeta-x)+O(q^3).
\]
It follows by direct projection that
\[
 d^{\rm self}=-e_\xi^\K\cdot b^\perp=1+O(q^2),
 \qquad c^{\rm self}=e_\xi^\K\cdot b=O(q).
\]
In particular the sign in the output-rotation channel is positive.
FWC \texttt{corefield} states, source by source,
\[
 y'_2=\frac{\rho}{2}(1-N)\kf(\rho x)+O(q^{1/10}),
 \qquad v_{\psi,2}=q y'_2,\qquad 2g_2=O(q^2).
\]
Therefore
\begin{align}
 q^{-2}O_{2,\K}^{\rm self}
 &=\frac{\rho}{2}(1-N)
       \int_\K(\zeta-x)\kf(\rho x)\dd\mu
                +O(q^{1/10}+q^2),\notag\\
 |P_{2,\K}^{\rm self}|&\le Cq^4.
 \label{an01bcv1:outputexpansion}
\end{align}
This calculation differentiates the exact source output formula,
not a static error remainder.

The function \((\zeta-x)\kf(\rho x)\) is uniformly bounded and
Lipschitz on the fixed chart. FWC's physical masses and distance,
or RTA's supremum tracking together with those masses, give
\[
 \left|\int_\K(\zeta-x)\kf(\rho x)\dd\mu
           -(\zeta-X_*)\kf(\rho X_*)\right|\le Cq^\chi,
 \quad |N-n_\rho|\le Cq^{1/8}.
\]
Combining this with \eqref{an01bcv1:r2self} and
\eqref{an01bcv1:outputexpansion} proves the first bound in
\eqref{an01bcv1:r2limit}. Since \(\gamma_*>\chi\),
\(L^2q^{\gamma_*}=o(q^\chi)\); also
\(q^{1/8},q^{1/10},q^2\le q^\chi\) for \(0<q\le1\).
This proves its stated explicit power.

For positivity, \(s\le b_J\), \(\lambda\le49/100\), and hence
\[
 1-n_\rho(s)=\frac1{1+e^{\lambda s}}\ge\delta_n.
\]
WIT gives \(\zeta-X_*\ge g_*/4\), and
\(\rho X_*\ge-7R/10\). The Gaussian function
\[
 \kf(t)=\phi(t)+t\Phi(t)=\int_0^\infty u\phi(u-t)\dd u
\]
is strictly positive and increasing at every finite t.
For a fully elementary positive bound,
\(\kf(-a)\ge\int_1^2u\phi(u+a)\dd u
\ge\tfrac32\phi(a+2)>0\) when \(a\ge0\).
Multiplying \(\rho/2\ge13/40\), the logistic slack, the gap,
and this Gaussian moment proves \eqref{an01bcv1:cout}.
The final lower bound follows from \eqref{an01bcv1:r2limit}.
No positivity of \(X_*\), and no extension of WIT's \(s\le0\)
region theorem, was used.
\end{proof}

\section{Exact failed inequality and the closed subset}
\label{an01bcv1:stop}
Step 2 requires, among its separate remainder estimates,
\begin{equation}
 |r_2|\le C_*q^{2+\kappa},\qquad \kappa=1/20
 \quad\hbox{(or any fixed \(\kappa>0\))}.
 \label{an01bcv1:failed}
\end{equation}
Theorem \ref{an01bcv1:obstructiontheorem} would force
\[
 \tfrac12c_{\rm out}\le C_*q^\kappa,
\]
which is impossible for a q-independent \(C_*\) as \(q\downarrow0\).
This is a false estimate, not an unpaid hypothesis or a missing
tail budget. Even the coherent limit retains this channel.
P \texttt{atomicrates} explicitly identifies its second-source
output rotation before weak saturation.

For clarity, the already available RTA rate transfer and the theorem
above yield the rigorously closed two-channel statement
\begin{align}
 q^{-2}(H^{(1)}+H^{(2)})
       &=(X_*-Y_*)\,x'_{2,*}+O(q^\chi),\notag\\
 q^{-2}D_2^\tau
       &=(X_*-Y_*)\,x'_{2,*}+\mathcal O_*+O(q^\chi),
 \label{an01bcv1:twoway}\\
 x'_{2,*}
       &=\frac{\lambda}{2}
       [n_\rho V_*-X_*+(1-n_\rho)Y_*]\Phi(\rho X_*).\notag
\end{align}
Indeed subtract \eqref{an01bcv1:r2limit} from RTA's convergence
to P's coherent source rate. This keeps the entire input channel,
not a mean-angle approximation to a producer contribution.
It does not establish producer signs or the exact-lag comparison.

The necessary change to a future contract is to retain
\(\mathcal O_*\), or the exact \(O_{2,\K}\), in the leading
second-source mechanism. One may then seek an \(O(q^{2+\chi})\)
bound for \(r_2-q^2\mathcal O_*\), which is already supplied here.
The unchanged \(r_2\) cannot have that bound. No such change is
silently made to P or to the requested factor definitions.

\paragraph{Constant order and physical units.}
Keep all inherited RTA/FWC constants first; fix a finite continuation
window covering J and its radius R. Then fix \(\delta_n,c_{\rm out}\)
and \(C_{\rm rem}\), and choose q last with
\(C_{\rm rem}q^{1/20}\le c_{\rm out}/2\).
All constants are uniform on the original \(\mathcal P_*\times J\).
The lower bound on the physical remainder is
\(\mu_1^2c_{\rm out}q^2/2\). None of these is a numerical premise.

\paragraph{Stop and handoff.}
Inventory is complete for the proposed sources; Step 2 is stopped
at \eqref{an01bcv1:failed}. The other individual remainder
evaluations, signed exact lag, exposure comparison, producer
attribution and dominance mechanism are not pursued after this
failure. There is no claimed full initialized B-core theorem.
RTA's Theorem A and persistence are unaffected.

For a revised deep-regime S6 task, keep both rotation channels at
leading order and then test the signed exact lag and rate-weighted
producer split. S7/S8 additionally need a reached canonical/bridge
state and lag-relative precision when the lag is a small cancellation.
Neither an absolute norm estimate nor the present positive
output-rotation result supplies those items. No prior source,
checkpoint or README is modified.
\end{document}
